% ---
% id: protocol-compiler/req:gate-outcomes
% kind: req
% title: "Outcome of each gate failure"
% summary: >-
%   A failing check or a failing certificate is refused with the check or the
%   condition named, whatever else fails; an unlicensed list or an absent
%   attack-set licence alone is emitted marked not key-length-eligible. The
%   engine never supplies a missing attestation or proceeds as though it had one.
% prominence: standard
% status: accepted
% uses: [thm:gates, def:eligible, def:gate-outcome-table, transmitter/req:licence-scope]
% mentions: [def:compile, transmitter/req:marginal-fill, transmitter/req:failing-certificate, transmitter/def:source-spec, foundations/req:consumer-decides-exactly]
% assumes: []
% refs: []
% provenance: newly-derived
% implements: certiqs_key.protocol_compiler.gates:resolve_gates
% lean:
% ---

The \tref{gate outcome table} (Definition~\ref{def:gate-outcome-table}) fixes which
of the two outcomes each failure has, as Theorem~\ref{thm:gates}\,(ii) states;
this requirement names the strings and the flag that carry them. A check that
fails, so that no declaration is shown to be well-formed, or a certificate
that does not verify, is a refusal; a \tref{well-formed declaration} on
verified data whose key-length claim is unsupported is emitted and marked. The
transmitter refuses a request whose certificate fails a condition of its form
(\xref{transmitter/req:failing-certificate}); gate (G2) re-decides the
conditions all the same, a consumer deciding the data it is handed
(\xref{foundations/req:consumer-decides-exactly}).

\begin{requirement}\label{req:gate-outcomes}
For each gate failure the engine produces the outcome the gate outcome table
(Definition~\ref{def:gate-outcome-table}) fixes for it, and no other; where
the failures of one declaration meet the premises of a refusal row and of an
emission row both, it produces the refusal:
\begin{enumerate}[label=(\roman*),leftmargin=2.2em]
\item row~1, a check that fails: a refusal naming the check, which in the
  general format covers \iface{ensemble_form_in_untrusted_mode},
  \iface{alice_ensemble_marginal_mismatch}, \iface{alice_ensemble_sum_rounds},
  \iface{alice_sector_marginal_zero} and
  \iface{constraint_scope_not_a_factor};
\item row~3, \tref{trusted-source mode} at any $d_A$ with a certificate present and
  failing a condition of gate (G2) for its form: a refusal naming the
  condition, an \tref{ensemble
  presentation} beside a POVM-form declaration being refused with
  \iface{ensemble_presentation_without_ensemble_form};
\item rows~5 and~6, a constraint list not in the licensed class or the
  attack-set \tref{licence record} \iface{attack_set_license} absent: an emission
  with \iface{key_length_eligible} false.
\end{enumerate}
\emph{A trusted-mode specification at $d_A\ge2$ whose certificate is absent
does not reach this module.} The mode is read off the transmitter artefact's
\tref{marginal slot} (Definition~\ref{def:compile}), and the transmitter
refuses a trusted-mode request at $d_A\ge2$ that carries no certificate before
any artefact is emitted, with \iface{trusted_mode_without_certificate}
(\xref{transmitter/req:licence-scope}; the modes are defined in
\xref{transmitter/def:source-spec}). So an artefact in trusted mode at
$d_A\ge2$ carries the record gate (G2) tests for, and the two rows of the
table that describe a certificate-absent case --- rows 2 and 4, at $m=0$ and
at $m\ge1$ --- are \emph{unreachable through the transmitter artefact}. At
$d_A=1$ the gate demands no certificate at all, the slot holds the $[1]$ the
transmitter filled it with (\xref{transmitter/req:marginal-fill}), and one
supplied there and failing is refused by row 3, its scope being every $d_A$.
Rows 2 and 4 are unreachable in the general format for the same reason, which
reads the transmitter's refusal and no property of the general format. The
engine never supplies a missing attestation and never proceeds as though it
had one.
\end{requirement}
